\chapter{Causality, Mechanisms, and Complexity}
\label{ch:causality}

\begin{chapterguide}
This chapter distinguishes association, intervention, and explanation. It
introduces potential outcomes, directed acyclic graphs, structural causal
models, interventionist explanation, mechanisms, and transportability. The
central claim is that causal knowledge is possible when the inquiry creates or
locates stable constraints on change, not merely patterns of co-occurrence.
\end{chapterguide}

Causal claims answer a practical question: what would happen if something were
changed? A correlation describes how variables vary together. A causal claim
says that a difference in one variable would produce a difference in another
under specified conditions. The distinction is not metaphysical decoration.
Medical treatment, economic policy, education, and climate intervention all
depend on it.

Causality is difficult because the same association can arise from direct
causation, common causes, selection, measurement error, feedback, or chance.
No statistical method can identify a causal interpretation without assumptions
about design, time, variables, and the process generating the data. But the
presence of assumptions does not make causal knowledge impossible. It means
that causal claims should expose the assumptions that make them testable.

\section{Potential outcomes and the missing counterfactual}

The potential-outcomes framework imagines that each unit has an outcome under
each possible treatment. Let \(Y_i(1)\) be the outcome for unit \(i\) if it
receives treatment and \(Y_i(0)\) the outcome if it does not. The individual
causal effect is

\[
  \tau_i = Y_i(1)-Y_i(0).
\]

For a given unit, we observe only one of these two potential outcomes. This is
the fundamental problem of causal inference \citep{rubin1974,holland1986}. A
randomized experiment solves part of the problem by making treatment assignment
independent of potential outcomes in expectation. The average difference
between treatment and control groups can then estimate an average causal
effect, provided the treatment and outcome are defined and measured properly.

Randomization is not a universal substitute for thought. It may have limited
external validity, noncompliance, attrition, spillovers, or ethical
constraints. A trial can estimate what happened under its treatment protocol
without establishing what would happen under a different dose, population, or
institution.

\section{Graphs and confounding}

A directed acyclic graph (DAG) represents assumed causal direction. In a
simple confounding structure, \(C\) affects both treatment \(T\) and outcome
\(Y\):

\begin{figure}[ht]
\centering
\begin{tikzpicture}[
  node distance=17mm,
  v/.style={draw=teal,fill=pale,circle,minimum size=8mm,font=\small},
  ar/.style={-{Latex[length=2mm]},thick,draw=teal}
]
\node[v] (c) {\(C\)};
\node[v,below left=of c] (t) {\(T\)};
\node[v,below right=of c] (y) {\(Y\)};
\draw[ar] (c) -- (t);
\draw[ar] (c) -- (y);
\draw[ar] (t) -- (y);
\end{tikzpicture}
\caption{A confounding graph: \(C\) is a common cause of treatment \(T\) and
outcome \(Y\). The association between \(T\) and \(Y\) can mix the treatment
effect with the path through \(C\).}
\label{fig:confounding}
\end{figure}

If \(C\) is measured and the graph is credible, conditioning on \(C\) may
block the back-door path from \(T\) to \(Y\). If \(C\) is unmeasured, researchers
may need randomization, an instrumental variable, a natural experiment,
negative controls, sensitivity analysis, or a different design. The graph
does not prove the causal story. It makes the story visible enough to argue
about.

Judea Pearl's structural causal models distinguish observation from
intervention. Informally, \(\Prob(Y\mid T=t)\) describes the outcome among
cases observed with \(T=t\), whereas \(\Prob(Y\mid \doop(T=t))\) describes what
would happen if \(T\) were set to \(t\) by an intervention that removes the
usual causes of treatment \citep{pearl2009}. The distinction explains why an
association can fail to answer a policy question.

\section{Bayesian reasoning about causal claims}

Bayesian reasoning can represent uncertainty about causal hypotheses. Suppose
\(H_1\) says a treatment has a positive effect and \(H_0\) says it has no effect.
The evidence \(E\) changes their odds according to

\[
  \frac{\Prob(H_1\mid E)}{\Prob(H_0\mid E)}
  =
  \frac{\Prob(E\mid H_1)}{\Prob(E\mid H_0)}
  \times
  \frac{\Prob(H_1)}{\Prob(H_0)}.
\]

The first ratio is the likelihood ratio: how much more expected the evidence
is under \(H_1\) than \(H_0\). The second ratio expresses prior odds. A strong
randomized design can make the likelihood ratio informative by reducing
confounding. A weak observational design may produce a large apparent effect
while leaving several causal models compatible with the data.

The equation does not remove causal assumptions. It shows where they enter. A
prior may encode background knowledge about mechanisms. The likelihood depends
on the measurement model, treatment assignment, missingness, and outcome
definition. Good causal inference therefore combines probability with design
and mechanism.

\section{Interventionist explanation}

James Woodward treats causal explanation in terms of stable relations under
interventions \citep{woodward2003}. A variable \(X\) explains a change in \(Y\)
when manipulating \(X\), while holding relevant background conditions fixed,
changes \(Y\) according to a relationship that is stable enough to support
counterfactual reasoning.

This account improves on correlation, but the phrase \enquote{holding fixed}
hides complexity. In a biological system, manipulating one gene may alter many
pathways. In an economy, changing interest rates can trigger expectations and
institutional responses. In a social setting, an intervention can change the
meaning of the variable itself.

\RCR\ therefore adds an \emph{intervention boundary}: a causal claim must state
which variables are changed, which processes are allowed to respond, and which
background conditions are treated as fixed. A causal law is not a universal
promise that every intervention works identically. It is a scope-bound
constraint on a family of interventions.

\section{Mechanisms and causal structure}

Mechanistic explanations describe parts, activities, organization, and the
way their interactions produce a phenomenon. A mechanism is not merely a list
of correlations. It explains how a change propagates through a system.

Mechanistic evidence is especially valuable when it connects different scales.
A drug may bind to a receptor, alter a signaling pathway, change cell behavior,
and improve a clinical outcome. Each link can fail independently. The
mechanistic story is stronger when it is supported by intervention and
measurement at several links rather than inferred from the final outcome alone.

Mechanisms are not always microscopic. A market institution, a school
organization, or a legal procedure can be a mechanism when its parts and rules
produce a stable social outcome. The relevant parts may be persons, norms,
resources, and sanctions rather than molecules.

\section{Complexity, feedback, and nonlinearity}

Complex systems create three challenges for causal inference.

First, effects can be nonlinear. A small change may have little effect below a
threshold and a large effect above it. A linear average can conceal this
structure.

Second, systems can contain feedback. Treatment changes behaviour; behaviour
changes treatment exposure; the outcome changes future measurement. A static
DAG may need to be expanded into a time-indexed graph.

Third, many causes can interact. The effect of an intervention may depend on
context, history, and network position. The average treatment effect can be
scientifically useful while failing to describe any particular subgroup.

These challenges do not defeat objectivity. They make \emph{transportability}
and \emph{heterogeneity} central. A result should be carried to a new setting
only after checking which mechanisms, distributions, and institutional
conditions are shared.

\begin{principlebox}
\textbf{Causal constraint principle.} A causal claim is objective to the degree
that it identifies a stable change-relation supported by design, measurement,
mechanism, or convergent intervention. The claim must state its intervention
boundary and the conditions under which the relation is expected to travel.
\end{principlebox}

\section{What causal knowledge does not require}

Causal knowledge does not require deterministic laws. An intervention can
change a probability distribution without fixing an individual outcome. Nor
does causal knowledge require a single metaphysical definition of cause.
Different sciences may use manipulation, process transmission, counterfactual
dependence, or mechanism construction as complementary modes.

Causal knowledge also does not require the elimination of all uncertainty. It
requires that the uncertainty be attached to a specified causal question. A
claim such as \enquote{this policy works} is too vague. A claim such as \enquote{for
schools with these resources, under this implementation, assignment to this
programme increases average attendance over one year by an estimated range} is
criticizable and therefore epistemically useful.

\section*{Chapter summary}

Correlation records association; causal inference concerns changes under
intervention. Potential outcomes clarify the missing counterfactual, DAGs make
confounding assumptions visible, and structural causal models distinguish
observation from intervention. Mechanisms explain how effects are produced.
Complexity requires attention to nonlinearity, feedback, heterogeneity, and
transportability. \RCR\ treats causal objectivity as the stabilization of
change-relations under explicit intervention boundaries.

\begin{discussion}
\begin{enumerate}
\item Can a randomized trial establish a causal effect if treatment changes
the outcome measure?
\item When is a mechanism more informative than a highly predictive model?
\item How should a causal claim be reported when its average effect hides
opposite effects in subgroups?
\end{enumerate}
\end{discussion}

\begin{furtherreading}
See Rubin (1974), Holland (1986), Pearl (2009), Woodward (2003), Hernan and
Robins (2020), Salmon (1984), and Craver (2007).
\end{furtherreading}

